\section{Where the five numbers come from}

\subsection{Two equations, and which is which}

There are two polynomials in play, and they are not the same polynomial. Mixing them up
is the easiest mistake to make here, so they are set side by side before anything else.

The first is the one this section is about. It is an ordinary least squares fit to the nine
GO-10 runs, with the two settings left in the units the bench works in --- current density
$J$ in $\unit{mA/cm^2}$ and feed ratio $R$ in $\unit{mol/mol}$ --- and with all five
estimable terms fitted at once:
%
\begin{equation}
  \widehat{\text{recovery}}\,[\%] \;=\; 33.0192 \;+\; 3.16999\,J \;-\; 14.4538\,R
  \;-\; 0.0954576\,J^{2} \;+\; 0.837907\,R^{2}
  \label{eq:real-units-fit}
\end{equation}

The second is the one \texttt{uv run limodel analyze} prints. The code does not fit in real
units and does not fit every term. It takes the logarithm of each setting, centres and scales
it to run from $-1$ to $+1$ across the runs, and then grows the model one term at a time,
keeping a term only while it lowers the leave-one-out error. For recovery it stops with three
terms and prints
%
\begin{equation}
  \widehat{\text{recovery}}\,[\%] \;=\; 11.5776 \;-\; 15.627\,x_2 \;+\; 6.08016\,x_2^{2},
  \qquad x_2 = \frac{\ln R - 1.395778}{1.039721}
  \label{eq:coded-fit}
\end{equation}

Equation~\eqref{eq:coded-fit} has no $J$ in it at all. Equation~\eqref{eq:real-units-fit} is
not what the code reports; it is the plain, all-terms, real-units fit, written out so that
every coefficient in it can be traced to arithmetic on nine rows of measurements. That is
what the rest of this section does. Why the code works in logged, coded units and selects its
terms is a separate matter and is not argued here.

\subsection{The question least squares asks}

Nine runs survived preprocessing, all on the GO-10 membrane. Each one is three numbers: the
current density it was held at, the feed ratio it was fed, and the lithium recovery that came
out of it.

\begin{table}[htbp]
  \centering
  \caption{The nine GO-10 runs the fit is made from.}
  \label{tab:nine-runs}
  \begin{tabular}{llrr}
    \toprule
    Run & $J$ ($\mathrm{mA/cm^2}$) & $R$ ($\mathrm{mol/mol}$) & Recovery (\%) \\
    \midrule
    AE & 13.6667 & 1.4277  & 38.27 \\
    AF & 13.6667 & 2.8554  & 18.34 \\
    AG & 13.6667 & 5.7108  & 6.16  \\
    AH & 7.5000  & 1.4277  & 36.59 \\
    AI & 7.5000  & 1.4277  & 37.66 \\
    AJ & 4.5000  & 1.4277  & 21.08 \\
    AN & 25.6667 & 1.4277  & 39.50 \\
    AP & 25.6667 & 1.4277  & 26.32 \\
    AT & 13.6667 & 11.4215 & 2.32  \\
    \bottomrule
  \end{tabular}
\end{table}

We choose a shape for the answer first: a second order polynomial in the two settings. Then
we ask for the five numbers $b_0,\dots,b_4$ that make the nine predicted recoveries as close
as they can be to the nine measured ones. ``As close as they can be'' has to be made precise,
because there is no single set of five numbers that hits all nine measurements. Least squares
makes it precise in one particular way: add up the squares of the nine misses, and pick the
five numbers that make that total as small as possible.
%
\begin{equation}
  b \;=\; \arg\min_{b_0,\dots,b_4} \;\sum_{i=1}^{9}
  \Bigl( y_i - \bigl[\, b_0 + b_1 J_i + b_2 R_i + b_3 J_i^{2} + b_4 R_i^{2} \,\bigr] \Bigr)^{2}
  \label{eq:least-squares-question}
\end{equation}

Here $y_i$ is the recovery measured on run $i$. Squaring, rather than taking the size of the
miss, is a choice, and it is worth knowing it was made: squaring makes a miss of $6$ count
four times a miss of $3$, so the answer is pulled hard towards the runs it would otherwise
fit worst. It also makes the problem solvable in closed form, which is the next subsection.

\subsection{The design matrix}

Equation~\eqref{eq:least-squares-question} is tidier in matrix form. Put the nine measured
recoveries in a column $y$, and build a matrix $X$ with one row per run and one column per
term. Each column is a rule applied to that run's own $J$ and $R$:

\begin{itemize}[nosep]
  \item column 1 is the constant $1$, the same for every run --- this is what $b_0$ multiplies;
  \item column 2 is $J$, the run's current density, as recorded;
  \item column 3 is $R$, the run's feed ratio, as recorded;
  \item column 4 is $J^2$, that same current density times itself;
  \item column 5 is $R^2$, that same feed ratio times itself.
\end{itemize}

Nothing is standardised, logged or re-centred. The numbers in $X$ are the bench's numbers and
their squares, which is exactly why the coefficients in \eqref{eq:real-units-fit} carry real
units.

{\footnotesize
\setlength{\arraycolsep}{4pt}
\begin{equation}
  X \;=\;
  \begin{array}{r}
    \text{AE}\\ \text{AF}\\ \text{AG}\\ \text{AH}\\ \text{AI}\\
    \text{AJ}\\ \text{AN}\\ \text{AP}\\ \text{AT}
  \end{array}\!\!
  \begin{bmatrix}
    1 & 13.6667 & 1.4277  & 186.7778 & 2.0383   \\
    1 & 13.6667 & 2.8554  & 186.7778 & 8.1532   \\
    1 & 13.6667 & 5.7108  & 186.7778 & 32.6128  \\
    1 & 7.5000  & 1.4277  & 56.2500  & 2.0383   \\
    1 & 7.5000  & 1.4277  & 56.2500  & 2.0383   \\
    1 & 4.5000  & 1.4277  & 20.2500  & 2.0383   \\
    1 & 25.6667 & 1.4277  & 658.7778 & 2.0383   \\
    1 & 25.6667 & 1.4277  & 658.7778 & 2.0383   \\
    1 & 13.6667 & 11.4215 & 186.7778 & 130.4511
  \end{bmatrix},
  \qquad
  y \;=\;
  \begin{bmatrix}
    38.27 \\ 18.34 \\ 6.16 \\ 36.59 \\ 37.66 \\ 21.08 \\ 39.50 \\ 26.32 \\ 2.32
  \end{bmatrix}
  \label{eq:design-matrix}
\end{equation}
}

Two features of this matrix decide everything that follows. Rows AH and AI are identical, and
so are rows AN and AP: those are repeats, two runs at the same setting. Identical rows must
get identical predictions, whatever the five numbers turn out to be. And the right-hand column
is nearly all the same value: six of the nine runs sat at $R = 1.4277$, so the feed ratio was
only ever moved at one current density.

\subsection{Why five columns and not six}

A full second order model in two settings has six terms, the sixth being the interaction
$J\cdot R$. It is left out of $X$ because over these nine runs it cannot be fitted at all --
not because it was judged unimportant.

The runs move current density only at the lowest feed ratio, and move feed ratio only at one
current density. That cross shape makes the product column an exact combination of three
columns already present:
%
\begin{equation}
  J_i R_i \;=\; 1.42769\,J_i \;+\; 13.66667\,R_i \;-\; 19.51176
  \qquad \text{for all nine runs}
  \label{eq:interaction-identity}
\end{equation}

This is not approximately true, it is true: the largest discrepancy across the nine rows is
$7.1\times10^{-15}$, which is floating point noise and nothing else. Add $J\cdot R$ as a sixth
column and the matrix has rank 5, not 6. Six unknowns, five independent pieces of information:
infinitely many different sets of six coefficients would fit the nine runs equally well, so
there is no answer to report. The code meets the same wall and handles it by refusing the
fit --- \texttt{SecondOrderPolynomial.find\_fit} checks the rank \texttt{numpy} hands back and
returns \texttt{None} when it falls short of the number of terms.

\subsection{The normal equations}

Minimising a sum of squares is a calculus problem. Write the total squared miss as
$S(b) = (y - Xb)^{\mathsf T}(y - Xb)$ and differentiate it with respect to each of the five
coefficients in turn. At the minimum every one of those five derivatives is zero, which gives
five equations in the five unknowns:
%
\begin{equation}
  \frac{\partial S}{\partial b} \;=\; -2X^{\mathsf T}(y - Xb) \;=\; 0
  \qquad\Longrightarrow\qquad
  \bigl(X^{\mathsf T}X\bigr)\,b \;=\; X^{\mathsf T}y
  \label{eq:normal-equations}
\end{equation}

These are the normal equations. The name comes from what they say: the vector of misses
$y - Xb$ is at right angles --- normal --- to every column of $X$. That has a plain reading.
If the leftover misses still leaned one way with $J$, you could lower the total by adjusting
$b_1$, so you were not at the minimum. At the minimum there is no leaning left to exploit.

Both sides are small, and both are made of nothing but sums of products of the numbers in
\eqref{eq:design-matrix}. The top left entry of $X^{\mathsf T}X$ is $9$, the number of runs,
because it is the first column dotted with itself. The second entry of the top row is
$125.5$, which is the nine current densities added up.

{\footnotesize
\setlength{\arraycolsep}{4pt}
\begin{equation}
  X^{\mathsf T}X =
  \begin{bmatrix}
    9        & 125.5      & 28.5538   & 2197.4167    & 183.4468   \\
    125.5    & 2197.4167  & 393.8044  & 44962.6528   & 2512.2024  \\
    28.5538  & 393.8044   & 183.4468  & 6070.4973    & 1716.9338  \\
    2197.4167& 44962.6528 & 6070.4973 & 1014258.2616 & 35316.4021 \\
    183.4468 & 2512.2024  & 1716.9338 & 35316.4021   & 18172.4790
  \end{bmatrix},
  \qquad
  X^{\mathsf T}y =
  \begin{bmatrix}
    226.2459 \\ 3230.6772 \\ 398.8249 \\ 60119.2880 \\ 1060.2725
  \end{bmatrix}
  \label{eq:normal-matrices}
\end{equation}
}

Note how wide the range of magnitudes is: $9$ in one corner and $1{,}014{,}258$ in another.
That is what squaring a current density of $25.67$ twice over does, and it matters in the next
subsection.

\subsection{Solving them}

Solving \eqref{eq:normal-equations} with those numbers gives the five coefficients:
%
\begin{equation}
  b \;=\;
  \begin{bmatrix} b_0 \\ b_1 \\ b_2 \\ b_3 \\ b_4 \end{bmatrix}
  \;=\;
  \begin{bmatrix}
    33.019212 \\ 3.169987 \\ -14.453829 \\ -0.0954576 \\ 0.837907
  \end{bmatrix}
  \label{eq:coefficients}
\end{equation}

which is equation~\eqref{eq:real-units-fit}, written out. That is the whole answer to ``how
those numbers'': they are what you get when you form the two arrays in
\eqref{eq:normal-matrices} out of nine rows of bench data and solve a five by five linear
system.

A solver does not take that route. \texttt{SecondOrderPolynomial.\allowbreak find\_fit}
calls \texttt{numpy.\allowbreak linalg.\allowbreak lstsq}, which asks the same question ---
it is still \eqref{eq:least-squares-question} --- but answers it by singular value
decomposition of $X$ itself, without ever forming $X^{\mathsf T}X$. Hand the five real-unit
columns of \eqref{eq:design-matrix} to either route and they agree to
$1.07\times10^{-12}$, so for reading the coefficients it makes no difference which is used.
(The code is never handed those columns: \texttt{find\_fit} only ever sees the logged, coded
columns of \eqref{eq:coded-fit}. What is shared is the solver, not the fit.)

A solver prefers the decomposition route because forming $X^{\mathsf T}X$ squares the
condition number, the measure of how much a solution can move when the inputs are nudged.
For this $X$ the condition number is $1770$; for $X^{\mathsf T}X$ it is
$3.13\times10^{6}$. Multiplying the matrix by its own transpose throws away about half the
available precision before the solving even starts, for no gain.

\begin{figure}[htbp]
  \centering
  \begin{tikzpicture}[
      box/.style={draw, rounded corners, align=center, inner sep=6pt, text width=7.4cm},
      lab/.style={font=\footnotesize, align=left, xshift=2pt}
    ]
    \node[box] (runs)
      {\textbf{Nine GO-10 runs}\\[2pt]
       each one a current density $J$, a feed ratio $R$\\ and one measured recovery};
    \node[box, below=0.95cm of runs] (design)
      {\textbf{Build the design matrix}\\[2pt]
       one row per run; columns $1,\;J,\;R,\;J^{2},\;R^{2}$\\
       (no $J\!\cdot\!R$: it is not independent here)};
    \node[box, below=0.95cm of design] (normal)
      {\textbf{Form the normal equations}\\[2pt]
       $\bigl(X^{\mathsf T}X\bigr)\,b = X^{\mathsf T}y$};
    \node[box, below=0.95cm of normal] (solve)
      {\textbf{Solve the $5\times5$ system}\\[2pt]
       by hand, or by any linear solver};
    \node[box, right=0.9cm of normal, text width=4.5cm] (svd)
      {\textbf{Or: skip $X^{\mathsf T}X$}\\[2pt]
       \texttt{numpy.linalg.lstsq},\\ by singular value\\ decomposition of $X$ itself};
    \node[box, below=0.95cm of solve] (out)
      {\textbf{Five coefficients}\\[2pt]
       $33.0192,\;3.16999,\;-14.4538,$\\ $-0.0954576,\;0.837907$};

    \draw[-Stealth] (runs)   -- node[right, lab] {$9 \times 3$ measurements} (design);
    \draw[-Stealth] (design) --
      node[right, lab] {$X$ is $9 \times 5$, \; $y$ is $9 \times 1$} (normal);
    \draw[-Stealth] (normal) --
      node[right, lab] {$X^{\mathsf T}X$ is $5 \times 5$, \; $X^{\mathsf T}y$ is $5 \times 1$}
      (solve);
    \draw[-Stealth] (solve) -- node[right, lab] {$b$ is $5 \times 1$} (out);
    \draw[-Stealth, dashed] (design.east) -| (svd.north);
    \draw[-Stealth, dashed] (svd.south) |-
      node[pos=0.86, below, lab, align=center, yshift=-1pt]
        {same $b$, to\\ $1.07\times10^{-12}$} (out.east);
  \end{tikzpicture}
  \caption{From nine runs to five coefficients. The solid path is the one this section
    works through: the shapes on its arrows are where the nine runs stop being nine things,
    since once $X^{\mathsf T}X$ and $X^{\mathsf T}y$ are formed the problem is five by five
    no matter how many runs went in. The dashed path is what a solver actually does --- the
    same least squares question, decomposing $X$ directly and never forming
    $X^{\mathsf T}X$. Both land on the same five numbers. Neither is the equation
    \texttt{uv run limodel analyze} prints: that one is fitted on the coded columns of
    \eqref{eq:coded-fit}.}
  \label{fig:least-squares-flow}
\end{figure}

\subsection{What the answer means run by run}

Put each run's $J$ and $R$ back into \eqref{eq:real-units-fit} and you get what the equation
thinks that run should have given. The difference is the residual.

\begin{table}[htbp]
  \centering
  \caption{Measured, fitted and residual recovery for all nine runs, with each run's
    leverage --- how much of its own fitted value the run itself dictates.}
  \label{tab:fitted-and-residuals}
  \begin{tabular}{lrrrr}
    \toprule
    Run & Measured (\%) & Fitted (\%) & Residual (\%) & Leverage \\
    \midrule
    AE & 38.27 & 39.59 & $-1.31$ & 0.644 \\
    AF & 18.34 & 24.07 & $-5.73$ & 0.337 \\
    AG & 6.16  & 3.30  & $+2.86$ & 0.834 \\
    AH & 36.59 & 32.50 & $+4.09$ & 0.258 \\
    AI & 37.66 & 32.50 & $+5.16$ & 0.258 \\
    AJ & 21.08 & 26.42 & $-5.34$ & 0.677 \\
    AN & 39.50 & 32.57 & $+6.93$ & 0.499 \\
    AP & 26.32 & 32.57 & $-6.25$ & 0.499 \\
    AT & 2.32  & 2.73  & $-0.41$ & 0.997 \\
    \midrule
    \multicolumn{3}{l}{Sum of squared residuals} & \multicolumn{2}{r}{$201.9127$} \\
    \bottomrule
  \end{tabular}
\end{table}

$201.9127$ is the smallest that total can be made with five numbers of this shape. Any other
five numbers give a larger total; that is what \eqref{eq:least-squares-question} asked for and
what \eqref{eq:coefficients} delivers.

The residuals are not scattered at random, and the leverage column says why. Leverage is the
share of its own fitted value that a run dictates; the nine leverages add to $5$, one for each
coefficient, so they are literally how the five coefficients are spent across the nine runs.

Run AT, the single run at $R = 11.4215$, has leverage $0.997$. It is almost the only evidence
the $R^2$ column has at the far end of the feed ratio range, so the surface bends through it,
and its residual is $-0.41$. AG, the only run at $R = 5.7108$, is the next most influential at
$0.834$ and is fitted to $+2.86$. In other words, where the design has one run and no second
opinion, the fit passes close to it almost by obligation, not because the physics was captured.

The big residuals sit at the opposite end. Rows AH and AI of $X$ are the same row, so both get
the same prediction, $32.50$; rows AN and AP are the same row, so both get $32.57$. But the
bench measured $36.59$ and $37.66$ at the first setting and $39.50$ and $26.32$ at the second.
The fit splits each pair down the middle and the disagreement between the repeats lands
entirely in the residuals: $+4.09$, $+5.16$, $+6.93$, $-6.25$. No choice of five coefficients
can remove those, because the four misses are the bench disagreeing with itself.

One run does not follow the simple reading. AF is the only run at $R = 2.8554$ yet it has the
\emph{lowest} leverage of the nine, $0.337$, and the second largest residual, $-5.73$. Being
alone at a setting is not what buys influence; sitting at the edge of the range is. AF sits in
the middle of the feed ratio range, where the quadratic in $R$ is already pinned down by the
runs on either side of it, so the fit is free to miss it.

Whether a total of $201.9127$ is small enough to believe is the next section's question.

\subsection{Why the coefficients cannot be read one at a time}

It is tempting to read $b_1 = +3.17$ as ``more current gives more recovery'' and
$b_4 = +0.838$ as ``more magnesium in the feed gives more recovery''. Both readings are wrong,
and they are wrong for the same reason: in a polynomial, a coefficient has no meaning on its
own. $J$ and $J^2$ move together, always, so $b_1$ and $b_3$ only say something as a pair.

Differentiate \eqref{eq:real-units-fit} with respect to $J$ and the pair speaks:
%
\begin{equation}
  \frac{\partial\,\widehat{\text{recovery}}}{\partial J}
  \;=\; 3.16999 \;-\; 2 \times 0.0954576\,J
  \;=\; 0 \quad\text{at}\quad J = 16.60\unit{mA/cm^2}
  \label{eq:j-turning-point}
\end{equation}

A positive $b_1$ with a negative $b_3$ is a hump, not a slope. The fitted surface rises with
current density up to about $16.6\unit{mA/cm^2}$ and falls after it, and that turning point
sits inside the range the runs covered, $4.5$ to $25.67$. The same pair of numbers read
separately would have said the opposite things about the two halves of that curve.

The feed ratio pair behaves the other way round. $b_2 = -14.45$ with $b_4 = +0.838$ turns at
$R = 8.62\unit{mol/mol}$, so over most of the tested range the surface falls steeply with
feed ratio and then flattens out as it approaches the highest run. The positive $b_4$ is not a
benefit of magnesium; it is the flattening of a steep fall.

The honest way to quote a fit like this is never a coefficient. It is a predicted recovery at
a stated setting.

\begin{table}[htbp]
  \centering
  \caption{What equation~\eqref{eq:real-units-fit} predicts at real settings. Reading down
    the first block, current density has a best value rather than a direction; reading down
    the second, feed ratio dominates everything.}
  \label{tab:predictions-at-real-settings}
  \begin{tabular}{rrr}
    \toprule
    $J$ ($\mathrm{mA/cm^2}$) & $R$ ($\mathrm{mol/mol}$) & Predicted recovery (\%) \\
    \midrule
    4.5000  & 1.4277 & 26.42 \\
    7.5000  & 1.4277 & 32.50 \\
    13.6667 & 1.4277 & 39.59 \\
    16.6042 & 1.4277 & 40.41 \\
    25.6667 & 1.4277 & 32.57 \\
    \midrule
    13.6667 & 2.8554  & 24.07 \\
    13.6667 & 5.7108  & 3.30  \\
    13.6667 & 11.4215 & 2.73  \\
    \bottomrule
  \end{tabular}
\end{table}

Table~\ref{tab:predictions-at-real-settings} says something no single coefficient could. Over
the whole tested span of current density the prediction moves by about
14 percentage points, and it moves up and then back down. Over the tested span of feed ratio,
at fixed current density, it moves by about 37 percentage points, all downward. The
five numbers in \eqref{eq:coefficients} encode that, but only together.
